\documentclass[12pt,a4paper,twoside,openright]{report}

% =========================================================
% LINGUA, CODIFICA, FONT
% =========================================================
\usepackage[utf8]{inputenc}
\usepackage[T1]{fontenc}
\usepackage[italian]{babel}

% Font stile Times/Times New Roman, compatibile con pdfLaTeX
\usepackage{newtxtext}
\usepackage{newtxmath}

% Migliora la resa tipografica e la giustificazione
\usepackage{microtype}

% =========================================================
% IMPAGINAZIONE
% =========================================================
\usepackage[
    a4paper,
    inner=3.5cm,   % margine interno maggiore per rilegatura
    outer=2.5cm,
    top=3cm,
    bottom=3cm,
    includeheadfoot
]{geometry}

% Interlinea: le linee guida indicano 1.15 o 1.5.
% Qui uso 1.5, più leggibile per una tesi.
\usepackage{setspace}
\onehalfspacing
% In alternativa, per interlinea 1.15:
% \setstretch{1.15}

% Rientro dei paragrafi
\setlength{\parindent}{1.25cm}
\setlength{\parskip}{0pt}
\usepackage{indentfirst}

% Pagine bianche senza intestazioni/numeri
\usepackage{emptypage}

% =========================================================
% NUMERAZIONE PAGINE
% =========================================================
\usepackage{fancyhdr}
\pagestyle{fancy}
\fancyhf{}

% Numero pagina esterno, utile per stampa fronte-retro
\fancyfoot[LE,RO]{\thepage}

\renewcommand{\headrulewidth}{0pt}
\renewcommand{\footrulewidth}{0pt}

\fancypagestyle{plain}{
    \fancyhf{}
    \fancyfoot[LE,RO]{\thepage}
    \renewcommand{\headrulewidth}{0pt}
    \renewcommand{\footrulewidth}{0pt}
}

% =========================================================
% TITOLI DEI CAPITOLI E SEZIONI
% =========================================================
\usepackage{titlesec}

% Capitoli nello stile: "1 Introduzione", senza scritta "Capitolo"
\titleformat{\chapter}[hang]
    {\normalfont\huge\bfseries}
    {\thechapter}
    {1em}
    {}

\titlespacing*{\chapter}{0pt}{-5pt}{25pt}

% Profondità numerazione e indice
\setcounter{secnumdepth}{3}
\setcounter{tocdepth}{2}

% =========================================================
% FIGURE, TABELLE, DIDASCALIE
% =========================================================
\usepackage{graphicx}
\graphicspath{{figures/}}

\usepackage{float}
\usepackage{subcaption}

\usepackage{caption}
\DeclareCaptionLabelSeparator{dot}{. }

% Didascalie simili alla tesi di esempio: Figura 1. ...
\captionsetup{
    font=small,
    labelfont=it,
    textfont=it,
    labelsep=dot
}

% Tabelle
\usepackage{booktabs}
\usepackage{array}
\usepackage{longtable}
\usepackage{multirow}
\usepackage{tabularx}
\usepackage{tikz}
\usetikzlibrary{arrows.meta,positioning,shapes.geometric,fit,calc}
\usepackage{pgfplots}
\pgfplotsset{compat=1.18}

% =========================================================
% MATEMATICA E UNITÀ DI MISURA
% =========================================================
\usepackage{amsmath}
\let\Bbbk\relax
\usepackage{amssymb}

% Equazioni numerate per capitolo: (2.1), (2.2), ...
\numberwithin{equation}{chapter}

% siunitx non è sempre presente nella distribuzione TeX installata localmente.
% Se vuoi usarlo, installa texlive-science oppure riattiva questo blocco.
\IfFileExists{siunitx.sty}{
    \usepackage{siunitx}
    \sisetup{
        locale=DE,
        detect-all,
        per-mode=symbol
    }
}{}

% =========================================================
% ELENCHI
% =========================================================
\usepackage{enumitem}
\setlist{
    noitemsep,
    topsep=4pt
}

% =========================================================
% CODICE
% =========================================================
\usepackage{xcolor}
\usepackage{listings}

\lstdefinestyle{tesiCode}{
    basicstyle=\ttfamily\small,
    breaklines=true,
    frame=single,
    numbers=left,
    numberstyle=\tiny,
    captionpos=b,
    tabsize=4,
    showstringspaces=false
}

\lstset{style=tesiCode}

\renewcommand{\lstlistingname}{Listato}
\renewcommand{\lstlistlistingname}{Elenco dei listati}

% =========================================================
% BIBLIOGRAFIA
% =========================================================
% Crea un file bibliografia.bib nella stessa cartella del main.tex
\usepackage{csquotes}
\usepackage[
    backend=biber,
    style=alphabetic,
    sorting=nyt
]{biblatex}

\addbibresource{bibliografia.bib}

% Per usare riferimenti stile [CabLZ16], sostituire:
% style=numeric
% con:
% style=alphabetic

% =========================================================
% LINK E RIFERIMENTI INCROCIATI
% =========================================================
\usepackage{xurl}

\usepackage[
    hidelinks
]{hyperref}

\usepackage[italian,nameinlink,noabbrev]{cleveref}

\crefname{chapter}{capitolo}{capitoli}
\Crefname{chapter}{Capitolo}{Capitoli}
\crefname{section}{paragrafo}{paragrafi}
\Crefname{section}{Paragrafo}{Paragrafi}
\crefname{figure}{figura}{figure}
\Crefname{figure}{Figura}{Figure}
\crefname{table}{tabella}{tabelle}
\Crefname{table}{Tabella}{Tabelle}
\crefname{equation}{equazione}{equazioni}
\Crefname{equation}{Equazione}{Equazioni}
\crefname{lstlisting}{listato}{listati}
\Crefname{lstlisting}{Listato}{Listati}

% =========================================================
% DATI FRONTESPIZIO
% =========================================================
\usepackage{etoolbox}

\newcommand{\TesiTitolo}{Controllo semaforico adattivo mediante strategie Max-Pressure: progettazione e valutazione sperimentale}
\newcommand{\TesiCandidato}{Jakub Pawel Homoncik}
\newcommand{\TesiRelatore}{Prof. Giacomo Cabri}

% Lascia vuoto se non hai correlatori
\newcommand{\TesiCorrelatori}{}
% Esempio:
% \newcommand{\TesiCorrelatori}{Prof.ssa Manuela Montangero}

\newcommand{\TesiAnnoAccademico}{2025--2026}

\hypersetup{
    pdftitle={\TesiTitolo},
    pdfauthor={\TesiCandidato},
    pdfsubject={Controllo semaforico adattivo},
    pdfkeywords={Sistemi di trasporto intelligenti, Controllo semaforico adattivo, Simulazione del traffico urbano, Max-Pressure, SUMO}
}

% Finche' le sezioni sono vuote, questo segnaposto crea un paragrafo
% invisibile e permette a LaTeX di interrompere correttamente la pagina.
% I richiami possono essere rimossi man mano che viene inserito il testo.
\newcommand{\tesischeletro}{\mbox{}\par}

% I capitoli successivi a Manhattan restano nel sorgente come traccia di
% lavoro, ma non vengono compilati finche' non sono stati redatti.
\newif\ifcapitolifuturi
\capitolifuturifalse

% =========================================================
% FRONTESPIZIO
% =========================================================
\newcommand{\frontespizio}{
    \begin{titlepage}
        \thispagestyle{empty}
        \centering

        {\Large\bfseries
        UNIVERSITÀ DEGLI STUDI DI\par
        MODENA E REGGIO EMILIA
        \par}

        \vspace{1cm}

        {\large
        Dipartimento di Scienze Fisiche, Informatiche e Matematiche
        \par}

        \vspace{0.5cm}

        {\large
        Corso di Laurea in Informatica
        \par}

        \vfill

        {\LARGE\bfseries
        \TesiTitolo
        \par}

        \vfill

        \begin{minipage}[t]{0.45\textwidth}
            \raggedright
            \textbf{Relatore:}\par\vspace{0.2cm}
            \TesiRelatore

            \ifdefempty{\TesiCorrelatori}
            {}
            {
                \par\vspace{0.8cm}
                \textbf{Correlatori:}\par\vspace{0.2cm}
                \TesiCorrelatori
            }
        \end{minipage}
        \hfill
        \begin{minipage}[t]{0.45\textwidth}
            \raggedleft
            \textbf{Candidato:}\par\vspace{0.2cm}
            \TesiCandidato
        \end{minipage}

        \vfill

        {\large
        Anno Accademico \TesiAnnoAccademico
        \par}

    \end{titlepage}

    \cleardoublepage
}

% =========================================================
% ESEMPIO DI USO DOPO \begin{document}
% =========================================================
\begin{document}

\frontespizio

\chapter*{Sommario}
\addcontentsline{toc}{chapter}{Sommario}

Il controllo semaforico a tempi fissi non utilizza informazioni sul traffico
osservato e può quindi assegnare verde a direzioni poco richieste mentre si
formano code su altri accessi. Questa tesi studia strategie di controllo
adattivo basate sul principio Max-Pressure, nel quale ogni incrocio seleziona
la fase da servire confrontando le code a monte e a valle dei movimenti
consentiti. Il lavoro comprende la progettazione di una piattaforma
sperimentale riproducibile basata su SUMO e TraCI, l'implementazione di un
controller Max-Pressure e di sette estensioni specialistiche, e la costruzione
di una pipeline di ablation capace di generare popolazioni, eseguire batch
paralleli e produrre automaticamente metriche e report.

La prima valutazione completa considera una rete sintetica Manhattan \(8
\times 8\), tre livelli di domanda e cinque seed. I 135 casi sperimentali sono
terminati senza errori e con censura nulla. Rispetto al programma semaforico
statico nativo, il miglior controller riduce il tempo medio di attesa del
\(64{,}26\%\) nel traffico leggero, del \(60{,}09\%\) nel traffico medio e del
\(53{,}05\%\) nel traffico elevato. I risultati mostrano inoltre che
l'aggiunta di una regola specialistica non produce necessariamente un
beneficio: alcune estensioni migliorano soltanto in specifici regimi, altre
rimangono inattive oppure introducono un costo superiore al vantaggio
ottenuto. Questa osservazione motiva la successiva applicazione alla rete
reale di Bologna e lo studio di controller capaci di adattare la propria
modalità operativa al regime locale di traffico.

\chapter*{Parole chiave}
\addcontentsline{toc}{chapter}{Parole chiave}

\begin{itemize}
    \item Sistemi di trasporto intelligenti;
    \item Controllo semaforico adattivo;
    \item Simulazione del traffico urbano;
    \item Max-Pressure;
    \item SUMO.
\end{itemize}

\cleardoublepage
\tableofcontents
\cleardoublepage

\chapter{Introduzione}
\label{chap:introduzione}

La regolazione semaforica determina come una risorsa limitata, il diritto di
passaggio attraverso un incrocio, viene distribuita tra flussi tra loro
conflittuali. Una scelta apparentemente locale può propagarsi lungo la rete:
un verde troppo breve lascia una coda residua, un verde inutile sottrae tempo
agli altri accessi e una coda che raggiunge l'incrocio precedente può impedire
movimenti altrimenti non conflittuali. Per questo motivo il controllo dei
semafori costituisce un problema centrale nei sistemi di trasporto
intelligenti.

Questo capitolo introduce il problema affrontato, gli obiettivi e i contributi
del lavoro. La trattazione rimane volutamente ad alto livello; le formulazioni
matematiche, le scelte implementative e il protocollo sperimentale vengono
sviluppati nei capitoli successivi.

\section{Contesto e motivazioni}

I piani semaforici a tempi fissi sono semplici, prevedibili e poco costosi da
gestire. Essi associano a ciascuna fase una durata definita in anticipo e
ripetono ciclicamente la stessa sequenza. La loro efficacia dipende tuttavia
dalla corrispondenza tra il profilo di domanda usato durante la progettazione e
il traffico effettivamente presente. Quando la domanda cambia nello spazio o
nel tempo, il piano può continuare a servire una direzione vuota mentre su un
altro accesso si accumulano veicoli.

Il controllo adattivo cerca invece di usare osservazioni correnti della rete
per modificare il servizio offerto. Tra le strategie decentralizzate, il
Max-Pressure è particolarmente interessante perché non richiede una previsione
completa della domanda né un unico ottimizzatore centrale. Ogni intersezione
confronta le code a monte e a valle dei movimenti consentiti e assegna il verde
alla fase con pressione maggiore. La letteratura ha mostrato che questa
famiglia di algoritmi possiede importanti proprietà di stabilità e throughput
in opportune ipotesi modellistiche \cite{varaiya2013maxpressure,wongpiromsarn2012distributed,levin2023review}.

Il passaggio dalla teoria a una simulazione microscopica realistica introduce
però costi e vincoli che il modello ideale tende a semplificare: giallo e
tempo di sgombero, minimo e massimo verde, capacità finita delle corsie,
spillback, arrivi a plotoni e necessità di non trascurare movimenti deboli. Ne
segue che ``scegliere sempre la pressione massima'' non è, da solo, una
specifica sufficiente per un controller utilizzabile.

\section{Problema affrontato}

Il lavoro affronta due problemi collegati. Il primo è ingegneristico: costruire
una piattaforma con cui implementare, eseguire e confrontare in modo
riproducibile controller semaforici diversi su reti SUMO sintetiche e reali. Il
secondo è sperimentale: stabilire in quali condizioni il Max-Pressure di base
o le sue estensioni migliorino effettivamente il flusso rispetto a un programma
statico.

L'indicatore principale è il tempo di attesa, affiancato da tempo di viaggio,
perdita di tempo, velocità, numero di viaggi completati e censura. Il lavoro
non persegue obiettivi ambientali: consumi di carburante ed emissioni non sono
utilizzati per selezionare o valutare i controller. Tale delimitazione mantiene
coerente l'analisi con la domanda di ricerca, cioè massimizzare l'efficienza
degli incroci e del flusso veicolare.

Una difficoltà metodologica importante è evitare conclusioni ottenute da un
solo seed o da una simulazione interrotta prima che tutti i veicoli completino
il viaggio. Un controller può mostrare una bassa attesa media semplicemente
perché i veicoli più problematici sono ancora nella rete. Per questo la
piattaforma registra esplicitamente la quota di viaggi non conclusi e considera
la censura una condizione necessaria per interpretare correttamente le altre
metriche.

\section{Obiettivi della tesi}

Gli obiettivi sono i seguenti:

\begin{enumerate}
    \item implementare un controller Max-Pressure decentralizzato compatibile
    con le logiche semaforiche presenti nelle mappe SUMO;
    \item introdurre estensioni mirate a switching loss, minimo verde dinamico,
    fairness, capacità downstream, spillback e arrivi a plotoni;
    \item realizzare una pipeline di ablation che generi popolazioni
    riproducibili, esegua casi in parallelo e produca report confrontabili;
    \item definire metriche capaci di distinguere un reale miglioramento da un
    risultato condizionato da viaggi incompleti;
    \item valutare separatamente le varianti su una rete Manhattan regolare,
    prima di affrontare la maggiore eterogeneità della rete di Bologna;
    \item ricavare dai risultati indicazioni progettuali per controller capaci
    di adattarsi a regimi di traffico differenti.
\end{enumerate}

\section{Contributi del lavoro}

Il progetto di partenza forniva il collegamento fondamentale con SUMO, un
runner, un generatore di popolazioni e un controller fisso. Su questa base sono
stati sviluppati:

\begin{itemize}
    \item un'implementazione completa del Max-Pressure, comprensiva di macchina
    a stati per preservare giallo e all-red;
    \item sette famiglie di estensioni parametriche e relativi contatori di
    attivazione;
    \item una rappresentazione uniforme degli scenari sperimentali e dei
    profili di guida;
    \item una pipeline batch con preflight, cache delle popolazioni,
    parallelismo, tracciamento del progresso e report automatici;
    \item il calcolo esplicito di throughput e censura a partire dai file
    \texttt{tripinfo};
    \item una campagna finale su Manhattan $8\times8$ composta da 135
    simulazioni valide, tre livelli di domanda e cinque seed;
    \item l'estensione successiva della metodologia a distribuzioni spaziali
    non uniformi e a controller ibridi per la rete reale di Bologna.
\end{itemize}

Il contributo non consiste quindi soltanto in un algoritmo, ma in un processo
sperimentale osservabile. I contatori interni consentono di verificare se una
variante ha prodotto un risultato perché il suo meccanismo si è attivato o se,
al contrario, è rimasta operativamente identica al Max-Pressure di base.

\section{Organizzazione della tesi}

Il \cref{chap:stato-arte} presenta il controllo semaforico adattivo, il
Max-Pressure e gli strumenti di simulazione. Il \cref{chap:progettazione}
definisce requisiti e architettura della piattaforma. Il
\cref{chap:implementazione} descrive l'implementazione del controller e delle
estensioni. Il \cref{chap:metodologia} formalizza il protocollo sperimentale,
le popolazioni e le metriche. Il \cref{chap:manhattan} analizza la campagna
Manhattan $8\times8$. La parte successiva della tesi applica le indicazioni
ricavate alla rete di Bologna, introduce i controller ibridi sviluppati per la
mappa reale e discute i risultati conclusivi.


\chapter{Contesto e stato dell'arte}
\label{chap:stato-arte}

Per collocare correttamente il contributo è necessario distinguere tre livelli:
le strategie di controllo già note in letteratura, gli strumenti software
utilizzati e le componenti effettivamente sviluppate durante il tirocinio. Il
capitolo procede dal controllo a tempi fissi al Max-Pressure, ne discute le
principali estensioni e conclude descrivendo lo stato iniziale del progetto.

\section{Controllo semaforico urbano}

Un'intersezione semaforizzata è descritta da un insieme di movimenti, ad
esempio attraversamento rettilineo o svolta, e da un insieme di fasi. Una fase
abilita contemporaneamente movimenti compatibili. Il controller deve decidere
quanto a lungo mantenere la fase corrente e quando avviare la transizione verso
la successiva. La qualità della decisione dipende sia dal criterio di scelta
sia dai vincoli fisici: un cambio di fase richiede intervalli durante i quali
la capacità utile si riduce.

\subsection{Controllo a tempi fissi}

Nel controllo a tempi fissi la sequenza, la durata delle fasi e l'eventuale
offset rispetto agli altri incroci sono definiti a priori. Il principale
vantaggio è la prevedibilità. Un piano ben calibrato può coordinare corridoi e
creare onde verdi; inoltre non dipende dall'affidabilità di sensori o
comunicazioni.

Il limite è la scarsa reattività. Una volta fissato il ciclo, il verde viene
assegnato anche in assenza di domanda e non può essere trasferito
immediatamente a una coda inattesa. Il confronto sperimentale deve tuttavia
essere formulato con cautela: battere un particolare piano statico non equivale
a dimostrare la superiorità rispetto a ogni possibile piano ottimizzato. Nella
campagna Manhattan il riferimento è il \emph{program0} nativo della mappa, non
un piano coordinato ottimizzato a posteriori sui casi di test.

\subsection{Controllo attuato e adattivo}

I controller attuati usano rivelatori per prolungare o terminare una fase sulla
base della presenza veicolare. Regole di \emph{gap-out}, ad esempio, chiudono il
verde quando l'intervallo tra passaggi supera una soglia
\cite{cesme2012gapout}. I sistemi adattivi operano a un livello più ampio:
stimano lo stato corrente e aggiornano durate, cicli o selezione delle fasi.
OPAC e SCOOT sono esempi storici di strategie sensibili alla domanda
\cite{gartner1983opac,robertson1991scoot}.

Le soluzioni centralizzate possono coordinare una rete, ma richiedono più
informazioni, comunicazione e capacità di calcolo. Le soluzioni decentralizzate
riducono tale dipendenza e risultano naturalmente scalabili. Il Max-Pressure
appartiene a quest'ultima famiglia: l'azione di ciascun semaforo dipende
principalmente dalle code sulle corsie incidenti.

\section{Strategie Max-Pressure}

Il Max-Pressure deriva dalle politiche di backpressure impiegate per
stabilizzare reti di code. Applicato al traffico, confronta la quantità di
veicoli che desidera attraversare l'incrocio con la capacità residua delle
corsie riceventi. L'intuizione è servire i movimenti che riducono maggiormente
il differenziale di coda senza trasferire indiscriminatamente congestione a
valle.

\subsection{Principio di funzionamento}

Sia $\mathcal{M}$ l'insieme dei movimenti di un semaforo. Un movimento
$m=(i,j)$ collega una corsia in ingresso $i$ a una corsia in uscita $j$. Nella
forma adottata nel progetto, la pressione del movimento è

\begin{equation}
    p_{ij}(t)=q_i(t)-q_j(t),
    \label{eq:movement-pressure}
\end{equation}

in cui $q_i$ e $q_j$ sono le code osservate all'istante $t$. Se
$\mathcal{M}_\phi$ è l'insieme dei movimenti abilitati dalla fase principale
$\phi$, la pressione della fase è

\begin{equation}
    P_\phi(t)=\sum_{(i,j)\in\mathcal{M}_\phi}p_{ij}(t).
    \label{eq:phase-pressure}
\end{equation}

Il candidato naturale è quindi

\begin{equation}
    \phi^*(t)=\arg\max_{\phi\in\Phi}P_\phi(t).
    \label{eq:mp-choice}
\end{equation}

La sottrazione della coda downstream distingue il Max-Pressure da una politica
che osserva soltanto gli ingressi: un movimento perde priorità se conduce a
una corsia già occupata. Formulazioni più generali includono turning ratio,
pesi o portate di saturazione. La scelta semplificata adottata è coerente con
l'obiettivo di ottenere un controller locale interpretabile e applicabile alle
mappe disponibili.

\subsection{Pressione dei movimenti e delle fasi}

La pressione non rappresenta la sola lunghezza della coda. Una fase può avere
molti veicoli in ingresso ma una pressione modesta quando anche le uscite sono
congestionate. Al contrario, una coda più corta può risultare prioritaria se
può essere scaricata verso corsie libere. Il punteggio di fase somma i
movimenti consentiti; ciò permette di valutare con un unico valore configurazioni
semaforiche che servono più approcci contemporaneamente.

In una rete regolare il calcolo è completamente locale e può essere ripetuto a
ogni passo. Le proprietà teoriche di throughput del Max-Pressure sono state
studiate in reti di intersezioni e in implementazioni decentralizzate
\cite{varaiya2013maxpressure,wongpiromsarn2012distributed,le2015decentralized}.
Studi di simulazione mostrano tuttavia che parametri, sensori e struttura della
rete influenzano le prestazioni osservate \cite{sun2018simulation,mercader2020practical}.

\subsection{Proprietà e limiti del controllo locale}

L'uso di informazioni locali rende l'algoritmo scalabile e robusto alla
mancanza di un coordinatore centrale. Una variazione di domanda modifica
immediatamente i differenziali di coda, senza richiedere la stima preventiva di
un intero profilo giornaliero.

Esistono però quattro limiti principali. Primo, il criterio ideale non
attribuisce automaticamente un costo al cambio di fase. Secondo, le code
misurate non descrivono sempre la capacità fisica residua di una corsia.
Terzo, massimizzare il servizio complessivo può ritardare ripetutamente un
movimento debole. Quarto, decisioni localmente corrette possono non produrre
coordinamento lungo un corridoio. La revisione di Levin evidenzia proprio la
distanza tra proprietà teoriche e dettagli necessari per implementazioni
pratiche \cite{levin2023review}.

\section{Estensioni del Max-Pressure}

Le varianti studiate nel progetto non sono controller indipendenti dal punto di
vista concettuale. Ognuna modifica il punteggio o il momento della decisione
per affrontare un limite specifico. Questo rende possibile una valutazione per
ablation: si mantiene il nucleo Max-Pressure e si osserva il contributo della
singola estensione.

\subsection{Costo delle transizioni semaforiche}

Durante giallo e all-red nessuna delle direzioni principali utilizza pienamente
l'intersezione. Cambi troppo frequenti possono quindi ridurre la capacità
nonostante ogni scelta sia localmente giustificata. Le strategie
\emph{lost-time-aware} introducono una soglia: la nuova fase deve offrire un
guadagno sufficiente a compensare il servizio perso durante la transizione.
La considerazione esplicita dello switching loss è un tema rilevante anche in
formulazioni apprese del Max-Pressure \cite{wang2022switching}.

Un minimo verde assolve una funzione simile, ma in modo temporale: impedisce di
rivalutare immediatamente la fase appena attivata. Un massimo verde costituisce
invece una salvaguardia contro permanenze eccessive e forza periodicamente una
nuova decisione.

\subsection{Prevenzione dello spillback}

Lo spillback si verifica quando la coda occupa la corsia fino a bloccare
l'intersezione a monte. In questa condizione il solo conteggio dei veicoli
fermi può reagire troppo tardi o non distinguere correttamente la capacità
residua. Le strategie capacity-aware usano occupazione, lunghezza della corsia
o stime di capacità per evitare di inviare veicoli verso un'uscita quasi piena
\cite{gregoire2015capacity,noaeen2021spillback}.

Nel progetto vengono valutate due forme. La penalità downstream riduce in modo
continuo il punteggio dei movimenti diretti verso corsie occupate. Il blocco
anti-spillback è invece una protezione discreta: oltre una soglia di
attivazione una fase viene esclusa, e torna disponibile soltanto sotto una
soglia di rilascio. L'isteresi evita oscillazioni quando l'occupazione si trova
vicino al limite.

\subsection{Fairness e prevenzione della starvation}

Una politica orientata al throughput può continuare a favorire flussi forti.
Per limitare la starvation, al punteggio di una fase può essere aggiunto un
bonus crescente con il tempo trascorso dall'ultimo servizio. Il bonus non
sostituisce la pressione, ma rende progressivamente più competitivo un
movimento trascurato. Il compromesso tra throughput e fairness è discusso
anche da Wu, Liu e Gettman \cite{wu2018fairness}.

\subsection{Gestione dei plotoni di veicoli}

I veicoli non raggiungono necessariamente l'incrocio in modo uniforme. Il
rilascio di un semaforo a monte può generare un plotone, cioè una sequenza di
passaggi con headway ridotto. Interrompere il verde a metà del plotone spreca
parte dell'opportunità di servizio e può ricreare rapidamente una coda. La
platoon extension osserva gli intervalli di passaggio vicino alla stop line e,
entro un limite massimo, prolunga la fase quando il flusso rimane compatto.
Questa logica è affine ai criteri di gap-out dei controller attuati, ma viene
integrata nella decisione Max-Pressure.

\section{Simulazione microscopica del traffico}

La valutazione di un controller richiede di rappresentare non soltanto flussi
aggregati, ma anche partenze, accelerazioni, cambi di corsia e interazioni tra
veicoli. Per questo il progetto impiega una simulazione microscopica.

\subsection{SUMO}

SUMO, \emph{Simulation of Urban MObility}, è un simulatore open source,
microscopico e a passi discreti, progettato per reti stradali di dimensione
urbana \cite{behrisch2011sumo,sumodocs2026}. Le reti sono descritte tramite file
XML contenenti archi, corsie, connessioni, incroci e programmi semaforici. La
simulazione produce informazioni per veicolo e consente di eseguire esperimenti
senza interfaccia grafica, requisito essenziale per campagne batch.

Nel progetto SUMO gestisce la dinamica veicolare e la validità delle transizioni
semaforiche. Il codice Python non sostituisce il simulatore: ne modifica le
azioni attraverso un'interfaccia di controllo e raccoglie le osservazioni
necessarie.

\subsection{Interfaccia TraCI}

TraCI, \emph{Traffic Control Interface}, espone un protocollo con cui un
processo esterno può interrogare e comandare una simulazione SUMO
\cite{wegener2008traci,tracidocs2026}. A ogni passo il controller legge, tra le
altre grandezze, veicoli fermi, occupazione delle corsie, stato delle fasi e
passaggi su rivelatori virtuali. Può quindi richiedere il passaggio a una fase
o a un programma semaforico differente.

Questa separazione è utile scientificamente: il modello del traffico rimane
quello di SUMO, mentre la politica di controllo è implementata in Python e può
essere modificata senza rigenerare la rete.

\section{Situazione iniziale del progetto}

Le linee guida di tesi richiedono di distinguere quanto già disponibile da
quanto realizzato nel lavoro. La ricostruzione è stata effettuata attraverso la
cronologia Git e il confronto tra il commit iniziale applicativo e la versione
finale.

\subsection{Componenti software già disponibili}

La base iniziale comprendeva:

\begin{itemize}
    \item un runner capace di avviare SUMO e avanzare la simulazione;
    \item la classe astratta comune ai controller;
    \item un controller a tempi fissi;
    \item un generatore di popolazioni YAML;
    \item strutture per veicoli, tipi e metriche elementari;
    \item utility per risolvere i percorsi delle mappe;
    \item una raccolta di reti sintetiche e reali, tra cui Manhattan e Bologna.
\end{itemize}

Era quindi già presente l'infrastruttura minima per inserire veicoli e
interagire con i semafori. Non erano invece disponibili un Max-Pressure
completo, le varianti specialistiche, una campagna automatizzata di ablation e
la verifica esplicita dei viaggi non conclusi.

\subsection{Limiti della soluzione iniziale}

L'esecuzione di una singola simulazione non è sufficiente per un confronto
sperimentale. Mancavano una matrice dichiarativa dei casi, isolamento degli
output, esecuzione parallela, controllo preventivo delle configurazioni,
aggregazione su seed e report delle differenze rispetto a una baseline. La
sola media sui veicoli arrivati, inoltre, poteva nascondere un'elevata censura.

Dal lato del controllo, un template intelligente non risolveva il problema
della mappatura tra segnali SUMO, corsie e movimenti, né la gestione sicura di
giallo e all-red. Questi aspetti sono risultati più complessi del calcolo della
sola pressione e hanno richiesto una macchina a stati esplicita.

\subsection{Contributi sviluppati nella tesi}

L'evoluzione del repository mostra una progressione incrementale: prima il
Max-Pressure puro, quindi spillback, lost-time-aware, fairness, Nmin, penalità
downstream e plotoni; successivamente la pipeline di ablation, la correzione
delle metriche e i controller ibridi per Bologna. La scelta di aggiungere una
funzione alla volta ha permesso di attribuire i risultati e di individuare
estensioni inattive o controproducenti. I capitoli seguenti descrivono questa
architettura e l'implementazione risultante.

\chapter{Requisiti e progettazione del sistema}
\label{chap:progettazione}

Un esperimento sui controller semaforici è credibile soltanto se la piattaforma
impedisce confronti non omogenei e conserva abbastanza informazioni per
riprodurre ogni risultato. Da questa esigenza derivano i requisiti e
l'architettura descritti nel capitolo. La progettazione separa la politica di
controllo dalla generazione del traffico e dalla valutazione, in modo che una
modifica a un componente non alteri implicitamente gli altri.

\section{Requisiti funzionali}

Il sistema deve permettere di:

\begin{enumerate}
    \item selezionare una mappa SUMO e una popolazione di veicoli;
    \item eseguire sia il programma semaforico nativo sia un controller
    Max-Pressure configurabile;
    \item associare a ogni variante un insieme esplicito di parametri;
    \item generare più livelli e distribuzioni di domanda usando seed noti;
    \item avanzare la simulazione fino all'esaurimento dei viaggi o a un limite
    dichiarato;
    \item salvare metriche per veicolo, metriche aggregate e contatori interni
    del controller;
    \item costruire automaticamente il prodotto cartesiano tra mappe, domande,
    seed e scenari;
    \item eseguire più casi in parallelo senza collisioni tra file;
    \item produrre classifiche e differenze rispetto a una baseline scelta.
\end{enumerate}

Un requisito ulteriore, emerso durante lo sviluppo, è la capacità di rilevare
errori prima di avviare campagne lunghe. La fase di preflight deve quindi
controllare nomi delle mappe, scenari, file di rete, opzioni dei controller e
percorsi di output.

\section{Requisiti non funzionali}

Oltre al risultato numerico, sono rilevanti riproducibilità, modularità e
osservabilità. Questi requisiti hanno guidato scelte quali la serializzazione
delle popolazioni, l'uso di configurazioni risolte e l'introduzione di contatori
per le regole specialistiche.

\subsection{Riproducibilità}

Due controller devono essere confrontati sugli stessi veicoli, con uguali
route, istanti di partenza e caratteristiche di guida. Per ogni coppia
mappa--domanda--seed viene quindi generato un unico file di popolazione,
riutilizzato da tutti gli scenari. Il seed compare sia nel nome sia nel
contenuto del file. La directory di una campagna conserva inoltre la
configurazione effettivamente risolta, evitando che un valore predefinito
modificato in seguito renda ambiguo un risultato precedente.

La riproducibilità non implica che seed diversi producano lo stesso valore;
serve piuttosto a rendere controllata la variabilità. Nella valutazione finale
Manhattan ogni scenario viene eseguito sui medesimi cinque seed, consentendo
confronti appaiati.

\subsection{Modularità}

Il simulatore, il controller, il generatore della domanda e la reportistica
hanno responsabilità distinte. Tutti i controller rispettano un'interfaccia
comune di collegamento e aggiornamento; il runner non deve conoscere la formula
specifica con cui viene scelta una fase. Analogamente, la pipeline batch
costruisce comandi per una simulazione singola senza duplicarne la logica.

Le estensioni del Max-Pressure sono abilitate da flag e parametri. Questo
consente di mantenere un solo nucleo implementativo e di definire scenari
sperimentali come configurazioni, anziché creare copie divergenti del
controller.

\subsection{Osservabilità}

Una differenza nelle metriche finali non rivela da sola il comportamento che
l'ha prodotta. Il sistema registra quindi numero di switch, attivazioni del
massimo verde, passi di hold Nmin, eventi di blocco e rilascio spillback,
prolungamenti per plotone e bonus di fairness. Se una variante coincide
numericamente con il controller base e il suo contatore rimane a zero, è
possibile concludere che il meccanismo non è stato sollecitato, non che abbia
fallito dopo essersi attivato.

Ogni worker scrive inoltre un log separato. Il file di progresso aggrega casi
completati, riusciti e falliti, rendendo possibile controllare una campagna
eseguita in background.

\section{Architettura generale}

La \cref{fig:architettura} mostra i componenti principali. La pipeline di
ablation genera la matrice dei casi e usa il runner per avviare processi SUMO
indipendenti. Il runner carica la popolazione, collega il controller mediante
TraCI e inoltra le osservazioni al modulo delle metriche. Al termine, gli
output dei casi vengono aggregati in report di campagna.

\begin{figure}[htbp]
    \centering
    \begin{tikzpicture}[
        node distance=8mm and 12mm,
        box/.style={draw, rounded corners, align=center, minimum height=11mm,
                    text width=34mm, fill=black!3},
        data/.style={draw, cylinder, shape border rotate=90, aspect=0.22,
                     align=center, minimum height=12mm, text width=29mm,
                     fill=black!3},
        arrow/.style={-{Latex[length=2.2mm]}, thick}
    ]
        \node[box] (ablation) {Pipeline di ablation\\\texttt{run\_ablation.py}};
        \node[box, right=of ablation] (runner) {Runner della simulazione\\\texttt{runner.py}};
        \node[box, right=of runner] (sumo) {SUMO\\dinamica microscopica};
        \node[data, below=of ablation] (population) {Popolazione YAML};
        \node[box, below=of runner] (controller) {Controller\\fixed o Max-Pressure};
        \node[box, below=of sumo] (traci) {TraCI\\osservazioni e comandi};
        \node[data, below=18mm of controller] (outputs) {CSV, tripinfo,\\contatori e log};
        \node[box, left=of outputs] (reports) {Aggregazione e report};

        \draw[arrow] (ablation) -- (runner);
        \draw[arrow] (population) -- (runner);
        \draw[arrow] (runner) -- (sumo);
        \draw[arrow] (sumo) -- (traci);
        \draw[arrow] (traci) -- (controller);
        \draw[arrow] (controller) -- (runner);
        \draw[arrow] (runner) -- (outputs);
        \draw[arrow] (outputs) -- (reports);
        \draw[arrow] (ablation) |- (reports);
    \end{tikzpicture}
    \caption{Architettura della piattaforma sperimentale.}
    \label{fig:architettura}
\end{figure}

\subsection{Runner delle simulazioni}

\texttt{runner.py} è il punto di ingresso per un singolo caso. Il modulo
interpreta gli argomenti, individua i file della mappa, avvia SUMO, carica la
popolazione e istanzia il controller richiesto. Il ciclo principale continua
finché \texttt{getMinExpectedNumber()} segnala veicoli presenti o ancora attesi,
oppure finché viene raggiunto il limite opzionale di passi.

A ogni iterazione il runner avanza SUMO e invoca l'aggiornamento del
controller. Alla chiusura legge il file \texttt{tripinfo}, combina le
informazioni di viaggio con i contatori online e scrive i risultati del caso.
Questa sequenza centralizza il ciclo di vita della simulazione e garantisce che
fixed e Max-Pressure differiscano soltanto nella politica semaforica.

\subsection{Generazione delle popolazioni}

\texttt{generate\_population.py} seleziona route valide dalla mappa e genera
veicoli con istante di partenza, tipo e parametri di guida. Il risultato è un
file YAML leggibile e riutilizzabile. \texttt{src/population.py} traduce poi
questa descrizione in comandi TraCI durante l'esecuzione.

La separazione tra generazione e simulazione impedisce che due scenari
consumino in ordine diverso lo stesso generatore casuale. Il traffico non viene
rigenerato per ogni controller: una data popolazione è un input immutabile del
confronto.

\subsection{Controller semaforici}

La classe base definisce il contratto comune. Il controller fixed può lasciare
in esecuzione un programma esistente o modificarne selettivamente il verde. Il
controller Max-Pressure costruisce, per ogni semaforo, una rappresentazione di
fasi principali, movimenti, corsie e stato temporale. Tutti i semafori della
mappa vengono collegati all'avvio e aggiornati in modo decentralizzato.

La configurazione del Max-Pressure è parametrica. Gli scenari non contengono
codice, ma combinazioni di flag quali \texttt{--lost-time-aware},
\texttt{--nmin-dynamic} o \texttt{--spillback-block}. Questa scelta ha reso
possibile correggere il nucleo una sola volta e confrontare varianti che
condividono la medesima gestione delle transizioni.

\subsection{Raccolta e aggregazione delle metriche}

Le metriche online descrivono lo stato e le decisioni del controller; quelle
post-simulazione provengono dai viaggi completati. Per ciascun caso vengono
salvati un CSV per veicolo, un riepilogo e i log. La pipeline combina poi i
casi in:

\begin{itemize}
    \item \texttt{run\_results.csv}, una riga per simulazione;
    \item \texttt{summary\_by\_group.csv}, statistiche aggregate per scenario;
    \item \texttt{summary\_vs\_base.csv}, differenze rispetto alla baseline;
    \item \texttt{summary.md}, tabella leggibile della campagna;
    \item \texttt{winners.txt} e \texttt{table.txt}, classifiche sintetiche;
    \item \texttt{config\_resolved.yaml}, configurazione effettiva.
\end{itemize}

Sebbene il formato generico possa contenere campi prodotti da SUMO per altre
analisi, la tesi aggrega e interpreta esclusivamente metriche di mobilità e
completezza dei viaggi.

\section{Flusso di esecuzione di una simulazione}

Il flusso completo è il seguente:

\begin{enumerate}
    \item la pipeline risolve mappe, scenari, livelli di domanda e seed;
    \item per ogni combinazione mappa--distribuzione--domanda--seed genera o
    recupera dalla cache una popolazione;
    \item un worker prepara una directory isolata e avvia il runner;
    \item il runner avvia un processo SUMO con i file della mappa;
    \item i veicoli vengono inseriti secondo il file YAML;
    \item il controller analizza i programmi semaforici e costruisce il proprio
    stato interno;
    \item a ogni passo SUMO aggiorna veicoli e corsie, quindi il controller
    osserva lo stato e può richiedere una nuova fase;
    \item al termine vengono chiusi TraCI e SUMO e viene analizzato il
    \texttt{tripinfo};
    \item il worker registra successo o fallimento e aggiorna il progresso;
    \item completati tutti i casi, i risultati vengono raggruppati per mappa,
    distribuzione, domanda e scenario.
\end{enumerate}

La generazione delle popolazioni precede l'esecuzione parallela. In questo modo
nessun worker modifica un input condiviso mentre altri casi sono in corso.

\section{Scelte progettuali}

\subsection{Separazione tra controller e simulatore}

Si è scelto di non incorporare la logica adattiva nei file XML dei programmi
semaforici. Il controller Python osserva e comanda SUMO mediante TraCI. Il
vantaggio è la facilità di sperimentazione e di strumentazione; lo svantaggio è
un costo di comunicazione a ogni passo. Per le dimensioni studiate, la maggiore
chiarezza e modificabilità prevalgono su tale costo.

La separazione consente inoltre di preservare le transizioni definite nella
mappa. Il controller sceglie una fase principale target, ma non sintetizza
arbitrariamente una stringa di segnali potenzialmente insicura.

\subsection{Configurazione mediante parametri}

Le varianti sono definite in uno \emph{scenario pack}. Ogni scenario associa
un nome stabile al controller e ai flag. La configurazione dichiarativa riduce
gli errori nei comandi lunghi, permette di riportare esattamente i parametri
nei risultati e rende semplice escludere combinazioni non pertinenti. Per la
campagna Manhattan sono state scelte una variante rappresentativa per ciascuna
famiglia, anziché tutte le combinazioni esplorate durante il tuning.

\subsection{Gestione delle mappe reali e sintetiche}

Le mappe sintetiche hanno topologia regolare e sono adatte a isolare il
comportamento algoritmico. Le mappe reali contengono numeri diversi di corsie,
fasi eterogenee, svolte, segmenti brevi e piani semaforici specifici. Il layer
di risoluzione dei percorsi e l'analisi dinamica dei programmi evitano di
codificare nel controller gli identificativi di una singola rete.

Non tutte le strategie sono tuttavia trasferibili senza adattamento. Un
programma statico specifico della mappa, come il program0 di Bologna, non ha un
equivalente semantico universale. Per questo i capitoli sperimentali separano
la valutazione generale su Manhattan dagli sviluppi specifici della rete
reale.

\chapter{Implementazione dei controller Max-Pressure}
\label{chap:implementazione}

Il calcolo della pressione occupa poche righe; la parte più delicata consiste
nel trasformare quel calcolo in azioni compatibili con programmi semaforici
eterogenei. Il capitolo descrive prima la rappresentazione dei semafori e la
macchina a stati, quindi le estensioni valutate. Le formule corrispondono alla
versione effettivamente utilizzata nella campagna Manhattan.

\section{Rappresentazione dei semafori}

All'avvio il controller interroga TraCI per ottenere tutti i semafori, le
logiche disponibili e i link controllati. Per ciascun semaforo costruisce una
struttura contenente numero e durata delle fasi, fasi principali, movimenti per
fase, tempi di attesa interni e stato della transizione. La costruzione non
presuppone identificativi specifici della mappa.

Quando più logiche sono disponibili viene scelta quella con il maggior numero
di fasi principali, preferendo il programma già attivo se appartiene ai
candidati migliori. Questa regola evita di imporre universalmente un
identificativo, ad esempio \texttt{1}, che su una mappa reale potrebbe avere un
significato differente.

\subsection{Identificazione delle fasi principali}

Una fase è considerata principale se la stringa di stato contiene almeno un
segnale verde, \texttt{g} o \texttt{G}, e non contiene segnali gialli. Le fasi
solo gialle o rosse sono trattate come transizioni. Questa distinzione è
necessaria perché il Max-Pressure deve scegliere tra configurazioni di servizio,
non tra gli stati intermedi richiesti per passare in sicurezza da una
configurazione all'altra.

La classificazione è ricavata dalla logica SUMO e non da una lista scritta a
mano. Ciò permette allo stesso controller di operare su incroci con numeri di
fasi e lunghezze delle stringhe di segnale diversi.

\subsection{Costruzione dei movimenti}

\texttt{getControlledLinks()} associa ogni posizione della stringa semaforica
a uno o più collegamenti tra corsia in ingresso e corsia in uscita. Per ogni
fase principale il controller visita le posizioni verdi e costruisce
l'insieme

\begin{equation}
    \mathcal{M}_\phi=\{(i,j)\mid \text{il movimento }i\rightarrow j
    \text{ è verde nella fase }\phi\}.
\end{equation}

Eventuali duplicati vengono eliminati. Il risultato è una relazione esplicita
tra una fase e le corsie sulle quali misurare code, occupazione e passaggi. Se
una fase non abilita alcun movimento valido non partecipa alla selezione.

\subsection{Conservazione delle transizioni di sicurezza}

Il controller non passa direttamente da una stringa verde a un'altra. Quando
seleziona un target, avanza alla fase successiva prevista dalla logica e
memorizza il target come \texttt{pending\_target}. Durante giallo e all-red
attende la durata configurata nel programma SUMO. Se la sequenza raggiunge una
fase principale diversa dal target, effettua il riallineamento soltanto al
termine della transizione.

Questa soluzione conserva le durate di sicurezza della mappa e separa due
decisioni: \emph{quale} fase servire e \emph{come} raggiungerla. Un meccanismo
di recupero interviene se il semaforo rimane per almeno 20 secondi in uno stato
non principale, condizione considerata anomala. Tale salvaguardia è stata
introdotta dopo aver osservato stalli in campagne iniziali.

\section{Controller Max-Pressure di base}

\subsection{Misura delle code}

La coda su una corsia è misurata mediante
\texttt{getLastStepHaltingNumber()}, cioè il numero di veicoli considerati
fermi nell'ultimo passo. La misura viene memorizzata in una cache valida per il
passo corrente, perché la stessa corsia può comparire in più movimenti. Si
indica con $q_l(t)$ il valore associato alla corsia $l$.

La scelta del numero di veicoli fermi, anziché del numero totale presente,
rende il punteggio sensibile alla congestione effettiva vicino
all'intersezione. Non rappresenta tuttavia in modo completo veicoli in
avvicinamento o capacità fisica residua; le estensioni downstream e platoon
sono state introdotte anche per compensare questa limitazione.

\subsection{Calcolo della pressione}

Per ogni movimento $(i,j)$ si applica l'\cref{eq:movement-pressure}. La
pressione della fase è la somma dell'\cref{eq:phase-pressure}. In assenza di
estensioni, l'algoritmo produce inoltre la domanda della fase

\begin{equation}
    D_\phi(t)=\sum_{(i,j)\in\mathcal{M}_\phi}q_i(t),
    \label{eq:phase-demand}
\end{equation}

utilizzata dal minimo verde dinamico e, nei controller successivi, per
classificare il regime di traffico.

Una pressione negativa è ammessa: indica che le uscite sono più cariche degli
ingressi. Se tutte le fasi valide hanno un punteggio finito, viene comunque
scelta la migliore in senso relativo. Nel caso anti-spillback una fase può
ricevere valore $-\infty$ e viene esclusa.

\subsection{Selezione della fase}

Siano $S_\phi$ i punteggi, uguali alla pressione salvo bonus o penalità. La
fase migliore cambia il semaforo soltanto se

\begin{equation}
    \phi^*\neq\phi_c
    \quad\land\quad
    S_{\phi^*}>S_{\phi_c}+M,
    \label{eq:switch-condition}
\end{equation}

in cui $\phi_c$ è la fase corrente e $M$ il margine di switch. Nel controller
base $M=0$. La disuguaglianza stretta evita cambi in caso di pareggio. Se la
condizione è soddisfatta, il controller registra lo switch, imposta il target
pendente e avvia la sequenza di transizione.

\subsection{Vincoli di minimo e massimo verde}

Il valore predefinito del minimo verde è 10 secondi. Prima di tale durata la
fase non viene rivalutata. Il massimo verde è 120 secondi e opera come safety
cap: se esiste una fase alternativa finita, il controller può forzare il cambio
quando il limite è raggiunto. Minimo e massimo non determinano quindi un ciclo
fisso; delimitano la frequenza delle decisioni adattive.

È importante distinguere questo comportamento da un piano statico. Il fixed
mantiene una durata prestabilita indipendentemente dalle code. Il Max-Pressure
mantiene la fase almeno per il minimo, poi può conservarla o cambiarla in base
ai punteggi osservati.

\section{Macchina a stati del controller}

Per ogni semaforo vengono mantenuti fase principale attiva, target pendente,
tempo di permanenza in stati non principali, hold Nmin, tempo di svuotamento e
prolungamento per plotone. Queste variabili rendono il controller una macchina
a stati, non una funzione priva di memoria.

\subsection{Fase attiva e target pendente}

Nello stato ordinario il semaforo si trova in una fase principale. Il
controller misura la rete e può restare oppure scegliere un target. Quando
sceglie un target diverso, entra nello stato di transizione e non prende nuove
decisioni di pressione fino al completamento della sequenza. In questo modo un
punteggio che cambia durante il giallo non interrompe o inverte la manovra.

Raggiunto il target, lo stato attivo viene aggiornato e i temporizzatori della
nuova fase vengono inizializzati. I tempi di attesa delle fasi non servite sono
mantenuti soltanto quando è abilitata la fairness.

\subsection{Gestione di giallo e all-red}

La durata persa nel lasciare una fase è calcolata sommando le durate delle fasi
non principali che la seguono fino alla prossima fase principale:

\begin{equation}
    L_{\phi_c}=\sum_{k\in\mathcal{T}(\phi_c)}d_k,
    \label{eq:lost-time}
\end{equation}

con $\mathcal{T}(\phi_c)$ insieme degli stati di transizione e $d_k$ relativa
durata. Il valore viene riutilizzato sia dal margine lost-time-aware sia dal
calcolo Nmin. Non viene quindi assunto un giallo uguale per tutti gli incroci.

\subsection{Recupero dalle transizioni anomale}

Due controlli evitano stalli. Se esiste un target pendente e una fase di
transizione termina, il controller avanza o si riallinea al target. Se invece
non esiste un target ma il semaforo rimane in una fase non principale oltre la
soglia di 20 secondi, viene ripristinata la prima fase principale valida. Il
recupero non è una politica di traffico: è una protezione operativa contro
stati incoerenti o logiche di rete non previste.

\section{Lost-Time-Aware}

La variante lost-time-aware richiede che il guadagno di punteggio superi una
stima del costo della transizione. Il margine complessivo è

\begin{equation}
    M=\varepsilon_{abs}+\varepsilon_{rel}|S_{\phi_c}|
      +\gamma\,s\,L_{\phi_c},
    \label{eq:lta-margin}
\end{equation}

in cui $\gamma$ è il guadagno lost-time, $s$ una portata di saturazione
espressa in veicoli al secondo e $L_{\phi_c}$ il tempo perso. La variante finale
Manhattan usa $\gamma=0{,}60$ e $s=0{,}50$.

L'effetto atteso è ridurre switch marginali. Un valore eccessivo può tuttavia
rendere il controller lento a reagire e prolungare fasi non più ottimali. Per
questo l'efficacia non può essere dedotta dalla sola plausibilità della regola:
deve essere misurata insieme al numero di switch e agli interventi del massimo
verde.

\section{Minimo verde dinamico Nmin}

Nmin stima quanti veicoli sia opportuno servire prima di abbandonare una fase.
Dato il costo di transizione dell'\cref{eq:lost-time}, calcola

\begin{align}
    N_{cost} &= \alpha\,s\,L_{\phi_c}, \\
    N_{target} &= \max\{N_{floor},N_{cost},gD_{\phi_c}\},
    \label{eq:nmin-target}
\end{align}

in cui $\alpha$ pesa il costo, $N_{floor}$ è un minimo di veicoli, $g$ il
coefficiente dipendente dalla domanda e $D_{\phi_c}$ la domanda della fase. Se
la domanda è positiva, $N_{target}$ viene limitato a non superarla. La durata
di hold è

\begin{equation}
    H_{\phi_c}=\max\left\{G_{min}^{N},\frac{N_{target}}{s}\right\}.
    \label{eq:nmin-hold}
\end{equation}

Se la fase si svuota prima del termine, un timer consente il rilascio anticipato
invece di attendere inutilmente. La configurazione studiata usa
$\alpha=0{,}80$, $N_{floor}=2$, $G_{min}^{N}=4$ secondi, $g=0{,}25$ e rilascio
dopo un secondo di domanda nulla.

Questa configurazione modifica due aspetti: abilita l'hold dinamico e riduce da
10 a 4 secondi il minimo verde di base. Nell'interpretare i risultati occorre
quindi evitare di attribuire automaticamente ogni guadagno alla formula Nmin.
I contatori mostrano, infatti, che nel traffico basso l'hold dinamico non viene
quasi mai attivato; il beneficio deriva soprattutto dalla maggiore rapidità di
rivalutazione.

\section{Fairness}

Per ogni fase non servita che presenta domanda viene incrementato un tempo di
attesa interno $w_\phi$. Quando la fase è attiva, o quando la sua domanda si
annulla, il valore viene azzerato. Il bonus è

\begin{equation}
    B_\phi(w_\phi)=\mu\frac{w_\phi}{w_\phi+w_{1/2}},
    \label{eq:fairness-bonus}
\end{equation}

ed è aggiunto alla pressione. $\mu$ è il bonus asintotico e $w_{1/2}$ il tempo
al quale viene raggiunta metà del bonus massimo. La configurazione Manhattan
usa $\mu=5$ e $w_{1/2}=20$ secondi.

La funzione cresce rapidamente all'inizio e poi satura. Non garantisce un
ordine rigido delle fasi, ma rende progressivamente più difficile ignorare un
movimento con domanda persistente. Il contatore registra quante volte è stato
applicato un bonus positivo e la somma complessiva dei bonus.

\section{Penalità downstream}

La pressione di base sottrae già la coda a valle. La variante downstream
aggiunge una penalità continua basata sull'occupazione filtrata con media
mobile esponenziale:

\begin{align}
    \bar{o}_j(t) &= \eta o_j(t)+(1-\eta)\bar{o}_j(t-1), \\
    p'_{ij}(t) &= q_i(t)-q_j(t)-\beta\bar{o}_j(t).
    \label{eq:downstream-penalty}
\end{align}

La campagna usa $\beta=0{,}4$ e $\eta=0{,}30$. Il filtro riduce la sensibilità
a picchi di un solo passo. La penalità è morbida: una corsia occupata rende il
movimento meno attraente, ma non lo esclude. Il rischio è penalizzare anche una
normale occupazione in reti lontane dalla saturazione, introducendo un costo
senza prevenire alcun blocco reale.

\section{Anti-spillback}

Il blocco hard usa due stime della saturazione downstream. La prima è
l'occupazione filtrata $\bar{o}_j$. La seconda approssima il riempimento della
capacità di stoccaggio:

\begin{equation}
    f_j=\min\left\{1,
    \frac{q_j}{\max(1,\ell_j/7{,}5)}\right\},
    \qquad z_j=\max\{\bar{o}_j,f_j\},
    \label{eq:spillback-score}
\end{equation}

in cui $\ell_j$ è la lunghezza della corsia in metri e 7,5 metri è lo spazio
medio di jam assunto per veicolo. Una corsia non bloccata entra nello stato di
blocco quando $z_j\geq\theta_{on}$ e sono presenti almeno $h_{min}$ veicoli
fermi; rimane bloccata finché $z_j\geq\theta_{off}$. Poiché
$\theta_{off}<\theta_{on}$, la regola possiede isteresi.

Nella configurazione \texttt{on90\_off80}, $\theta_{on}=0{,}90$,
$\theta_{off}=0{,}80$ e $h_{min}=3$. Se anche un solo movimento di una fase
conduce a una corsia bloccata, l'intera fase riceve pressione $-\infty$. È una
scelta conservativa: protegge la rete dal trasferimento di congestione, ma può
sottrarre servizio anche agli altri movimenti della stessa fase. I contatori
di blocco, rilascio e passi bloccati servono a verificare se la protezione sia
mai entrata in funzione.

\section{Platoon extension}

Per ogni corsia in ingresso alla fase attiva viene posto un rivelatore virtuale
a 8 metri dalla fine. Il controller registra gli istanti in cui i veicoli
attraversano tale posizione e calcola la media degli ultimi quattro headway.
Un plotone è riconosciuto se la media non supera la soglia configurata e
l'ultimo passaggio non è più vecchio del tempo di gap-out.

Quando il controller starebbe per rivalutare la fase, la presenza del plotone
può prolungare il verde fino a un massimo aggiuntivo. La configurazione finale
usa soglia e gap-out di 2,1 secondi e un massimo di 4 secondi. Un guard
downstream impedisce l'estensione solo quando sono contemporaneamente abilitate
le misure di spillback o penalità downstream. Nello scenario standalone
\texttt{mp\_platoon\_x4} tali moduli non sono attivi; di conseguenza il
parametro di guardia dichiarato non modifica il comportamento. Questo dettaglio
implementativo è rilevante per non sovrainterpretare il test.

\section{Margini e stabilità dello switch}

Oltre al termine lost-time, il controller supporta un margine assoluto e uno
relativo. Con la sola componente relativa,

\begin{equation}
    M=\varepsilon_{rel}|S_{\phi_c}|.
\end{equation}

Lo scenario Manhattan \texttt{mp\_switch\_rel005} usa
$\varepsilon_{rel}=0{,}05$. L'obiettivo è introdurre una banda morta: una fase
alternativa deve migliorare il punteggio corrente di almeno il 5\%. Anche qui
il vantaggio potenziale, meno oscillazioni, si contrappone al rischio di
ritardare una decisione utile.

\section{Contatori di attivazione}

Il controller espone un dizionario di statistiche aggregato nel CSV della run.
Per i test dei primi sei capitoli vengono utilizzati:

\begin{itemize}
    \item numero di switch decisi dal margine e numero di switch forzati dal
    massimo verde;
    \item passi in cui Nmin ha mantenuto la fase;
    \item eventi di blocco, rilascio e passi con spillback attivo;
    \item passi di estensione dovuti a un plotone;
    \item numero e somma dei bonus di fairness positivi.
\end{itemize}

Queste misure non sono obiettivi da massimizzare. Un numero elevato di
attivazioni può indicare che la regola affronta un fenomeno frequente, ma anche
che interviene troppo aggressivamente. La loro funzione è spiegare le metriche
di mobilità e verificare la coerenza tra configurazione e comportamento.

\chapter{Metodologia e piattaforma sperimentale}
\label{chap:metodologia}

La qualità di una politica semaforica non può essere stabilita da una singola
traccia. Il protocollo deve controllare domanda, seed, durata e completezza dei
viaggi, oltre a rendere confrontabili gli output. Il capitolo definisce la
pipeline di ablation, i modelli di domanda, le metriche e le minacce alla
validità. La configurazione specifica di Manhattan è presentata nel
\cref{chap:manhattan}.

\section{Obiettivi della valutazione}

La valutazione risponde a quattro domande:

\begin{enumerate}
    \item il Max-Pressure riduce attesa e tempo di viaggio rispetto al programma
    statico nativo?
    \item il vantaggio rimane osservabile al crescere della domanda?
    \item quali estensioni migliorano il controller base e in quale regime?
    \item il comportamento interno conferma che la regola specialistica si è
    effettivamente attivata?
\end{enumerate}

Per rispondere, si adotta un confronto per ablation: ogni scenario abilita una
sola famiglia di modifiche rispetto a \texttt{mp\_base}. Le popolazioni sono
identiche tra controller e cambiano soltanto tra seed o livello di domanda. Il
piano fisso \texttt{fixed\_program0} funge da riferimento esterno, mentre
\texttt{mp\_base} permette di misurare il contributo marginale delle
estensioni.

\section{Pipeline di ablation}

\texttt{utils/run\_ablation.py} automatizza la campagna. Dati insiemi di mappe,
domande, distribuzioni, seed e scenari, costruisce tutti i casi, prepara gli
input, esegue i processi e aggrega gli output. Questa automazione è stata
necessaria perché comandi manuali lunghi avevano prodotto, nelle prime prove,
errori difficili da diagnosticare e risultati non sempre omogenei.

\subsection{Generazione dei casi sperimentali}

Il numero totale di casi è

\begin{equation}
    N_{case}=N_{mappe}\,N_{set}\,N_{domande}\,N_{seed}\,N_{scenari}.
    \label{eq:case-count}
\end{equation}

Prima dell'avvio, la pipeline risolve i nomi delle mappe e degli scenari, calcola
i preset di domanda effettivi e scrive la configurazione. Per ciascuna
combinazione di mappa, set, domanda e seed genera una sola popolazione. Tutti i
controller riferiti a quella combinazione ricevono lo stesso file.

Gli scenari sono raccolti in pack versionati nel codice. Il pack contiene nomi
interpretabili e flag completi; il report non deve quindi inferire a posteriori
quali parametri fossero stati usati.

\subsection{Esecuzione parallela}

I casi sono indipendenti e vengono assegnati a un pool di worker. Ogni worker
avvia un processo SUMO e un processo Python di controllo. Il parametro
\texttt{--jobs} limita il numero di casi simultanei: aumentarne il valore
riduce il tempo totale finché CPU, memoria e input/output non diventano il
collo di bottiglia.

Il parallelismo è tra simulazioni, non all'interno del Max-Pressure. Questa
scelta è adatta a una matrice di molti casi e mantiene deterministico il
controller di ciascuna simulazione. Le directory isolate impediscono che due
worker sovrascrivano \texttt{tripinfo}, log o metriche.

\subsection{Organizzazione delle run}

Ogni campagna riceve un identificativo del tipo
\texttt{run\_NNNN\_AAAAMMGG\_hhmmss}. La directory contiene popolazioni,
risultati aggregati e una sottodirectory per ciascun caso. Il file di progresso
riporta il totale pianificato e il conteggio di successi e fallimenti.

Questa organizzazione preserva la provenienza dei risultati e permette di
confrontare campagne successive senza mescolarne i file. I report finali usati
nel \cref{chap:manhattan} provengono interamente da una sola run completa.

\subsection{Report automatici}

Per ogni gruppo mappa--set--domanda--scenario vengono calcolate la media sui
seed e la deviazione standard campionaria. Se $x_s$ è la metrica del seed $s$,

\begin{align}
    \bar{x} &= \frac{1}{n}\sum_{s=1}^{n}x_s, \\
    s_x &= \sqrt{\frac{1}{n-1}\sum_{s=1}^{n}(x_s-\bar{x})^2}.
    \label{eq:mean-std}
\end{align}

La differenza percentuale rispetto alla baseline $b$ è

\begin{equation}
    \Delta_x=100\frac{\bar{x}-\bar{x}_b}{\bar{x}_b}.
    \label{eq:relative-delta}
\end{equation}

Per attesa e tempo di viaggio un valore negativo indica un miglioramento. Le
classifiche usano come criterio primario il tempo medio di attesa, ma la
conclusione viene controllata sulle altre metriche e sulla censura.

Quando si confrontano due controller sugli stessi seed, si considerano anche
le differenze appaiate $d_s=x_{b,s}-x_{c,s}$. L'intervallo di confidenza al
95\% è

\begin{equation}
    \bar{d}\pm t_{0{,}975,n-1}\frac{s_d}{\sqrt{n}}.
    \label{eq:paired-ci}
\end{equation}

Con soli cinque seed l'intervallo non sostituisce una campagna statistica più
ampia, ma aiuta a distinguere miglioramenti ampi e stabili da differenze di
pochi decimi di secondo.

\section{Generazione della domanda di traffico}

Una popolazione contiene identificativo, route, istante di partenza e tipo di
ogni veicolo. La finestra di generazione è distinta dalla durata totale della
simulazione: dopo l'ultima partenza la simulazione prosegue finché tutti i
veicoli completano il viaggio, salvo un limite esplicito.

I tipi sono estratti con probabilità 75\% autovetture, 10\% veicoli per
consegne, 10\% motocicli, 3\% camion e 2\% autobus. L'eterogeneità influenza
lunghezza, accelerazione, distanza minima e dinamica di partenza, rendendo la
popolazione meno artificiale di un insieme composto da sole autovetture.

\subsection{Livelli low, medium e high}

I nomi low, medium e high non corrispondono a valori universali. La pipeline
può applicare override per mappe note o una regola euristica basata sul numero
di route. Per reti con oltre 650 route la domanda media è 4000 veicoli; low è
il 60\% e high il 140\%, cioè rispettivamente 2400 e 5600 veicoli. Tutte le
partenze sono collocate nell'intervallo $[0,3600]$ secondi.

Questa calibrazione rende i livelli proporzionati alla varietà di percorsi
della mappa. Non garantisce però che ``high'' produca saturazione o spillback:
la gravità effettiva dipende anche da capacità, lunghezza delle route e
concentrazione spaziale. Nel capitolo dei risultati tale distinzione risulta
essenziale.

\subsection{Distribuzione balanced}

Nel set balanced ogni route ha uguale probabilità di essere scelta e gli
istanti di partenza sono uniformi nella finestra. È la configurazione usata per
la campagna Manhattan: distribuisce la domanda sull'intera griglia e permette
di osservare il comportamento generale del controller senza introdurre
corridoi caldi deliberati.

Balanced non significa che ogni arco riceva esattamente lo stesso flusso. Le
route hanno lunghezze e sovrapposizioni differenti; significa che il generatore
non applica pesi aggiuntivi alla loro selezione.

\subsection{Distribuzione skewed}

Nel set skewed la probabilità di una route è proporzionale a

\begin{equation}
    w_r=E_r^{1{,}4},
    \label{eq:route-weight}
\end{equation}

in cui $E_r$ è il numero di archi della route. Le partenze rimangono uniformi.
Le route più lunghe vengono quindi campionate più spesso e, condividendo più
segmenti, tendono a creare assi maggiormente caricati.

La scelta delle singole route non è fissata manualmente a ogni esecuzione: è
casuale ma riproducibile tramite seed. Si controlla la distribuzione, non la
lista esatta dei percorsi più trafficati. Questo set è stato introdotto per le
campagne sulla rete reale e non compare nel confronto Manhattan del
\cref{chap:manhattan}.

\subsection{Distribuzione peak}

Peak usa lo stesso campionamento spaziale di skewed e concentra anche le
partenze. Il generatore estrae da due distribuzioni beta ampie, una anticipata
e una posticipata nella finestra, quindi miscela il valore con un'estrazione
uniforme:

\begin{equation}
    u_{peak}=0{,}80\,u_{\beta}+0{,}20\,u_{uniforme}.
    \label{eq:peak-mixture}
\end{equation}

La scelta della prima componente avviene con probabilità 0,55. Il risultato
simula due ondate di domanda senza concentrare tutti i veicoli nello stesso
istante. Anche questa distribuzione è riproducibile e viene usata nella parte
successiva della tesi dedicata a Bologna.

\subsection{Seed e riproducibilità}

Un unico oggetto pseudo-casuale, inizializzato dal seed, determina partenze,
route e tipi. La popolazione viene poi ordinata per istante di partenza e
serializzata. Nella campagna Manhattan si usano i seed da 1 a 5. Poiché ogni
file è condiviso da tutti gli scenari, le differenze tra controller sono
appaiate e non dipendono da un diverso campione di veicoli.

\section{Profilo di guida human-light}

Il profilo \texttt{human\_light} riduce l'idealizzazione del comportamento
senza introdurre un tuning aggressivo. La velocità iniziale usa il valore
\texttt{desired}, non sempre il massimo consentito. Per ciascun tipo vengono
specificati tempo di reazione $\tau$, imperfezione di guida $\sigma$, intervallo
di azione, gap alle precedenze e ritardo di partenza.

\begin{table}[htbp]
    \centering
    \caption{Parametri principali del profilo \texttt{human\_light}.}
    \label{tab:human-light}
    \begin{tabular}{lccccc}
        \toprule
        Tipo & $\tau$ [s] & $\sigma$ & Azione [s] & Gap [s] & Startup [s] \\
        \midrule
        Autovettura & 1,45 & 0,55 & 1,0 & 1,40 & 0,80 \\
        Consegna    & 1,55 & 0,50 & 1,0 & 1,45 & 1,00 \\
        Camion      & 1,75 & 0,45 & 1,0 & 1,60 & 1,20 \\
        Autobus     & 1,75 & 0,45 & 1,0 & 1,60 & 1,20 \\
        Motociclo   & 1,25 & 0,65 & 1,0 & 1,25 & 0,45 \\
        \bottomrule
    \end{tabular}
\end{table}

Il profilo è identico per tutti i controller. Non cerca di riprodurre una
specifica popolazione cittadina, ma evita partenze istantanee e reazioni
perfette che potrebbero favorire irrealisticamente switch molto frequenti.

\section{Controller di confronto}

\subsection{Fixed program0}

\texttt{fixed\_program0} seleziona esplicitamente il programma con identificativo
\texttt{0} e lascia a SUMO l'esecuzione ciclica delle durate definite nella
mappa. Non osserva code e non modifica il piano. È la baseline principale di
Manhattan perché rappresenta il funzionamento semaforico nativo disponibile
senza controllo adattivo.

\subsection{Fixed tuned}

La piattaforma supporta anche \texttt{fixed\_tuned}, che forza a 30 secondi le
fasi principali del program0 mantenendo le transizioni. Questo baseline è
stato utile nelle campagne esplorative su Bologna, ma non viene incluso nella
campagna Manhattan finale: il suo parametro non deriva da un'ottimizzazione
specifica della griglia e avrebbe aggiunto un confronto non necessario al
primo studio per ablation.

\subsection{Max-Pressure di base}

\texttt{mp\_base} usa pressione upstream--downstream, minimo verde di 10
secondi, massimo verde di 120 secondi e margine nullo. Nessuna estensione è
abilitata. Esso ha due ruoli: controller adattivo generale e riferimento
interno per stabilire se una modifica specialistica aggiunga valore.

\section{Metriche}

\subsection{Tempo di attesa}

Per ciascun veicolo SUMO fornisce il tempo totale trascorso in attesa. La media
è l'indicatore primario perché misura direttamente la conseguenza delle scelte
semaforiche. Viene riportato anche il 95-esimo percentile, utile a verificare
che una buona media non nasconda una coda di utenti fortemente penalizzati.

\subsection{Tempo di viaggio e time loss}

Il tempo di viaggio è la differenza tra arrivo e partenza. La \emph{time loss}
è la perdita rispetto a una percorrenza ideale secondo il modello SUMO e
comprende rallentamenti, arresti e interazioni. Le due metriche non sono
intercambiabili: route più lunghe aumentano il viaggio anche a parità di
controllo, mentre il time loss isola meglio l'inefficienza incontrata.

\subsection{Velocità media}

Per ogni viaggio completato la velocità media è il rapporto tra distanza e
durata; il report aggrega poi i veicoli. Un controller efficace dovrebbe
produrre, a parità di route, velocità maggiori coerentemente con minori attese
e perdite di tempo.

\subsection{Viaggi completati e throughput}

Il throughput osservato è il numero di viaggi completati. Quando la simulazione
prosegue fino a svuotamento, tutti gli scenari dovrebbero completare la stessa
popolazione. In presenza di un limite temporale, il numero di arrivi diventa
invece una metrica essenziale: un controller che lascia veicoli bloccati non
può essere valutato soltanto sui pochi arrivati.

\subsection{Censura}

Siano $N_p$ i viaggi pianificati e $N_c$ quelli completati. La censura è

\begin{equation}
    C=\frac{N_p-N_c}{N_p}.
    \label{eq:censoring}
\end{equation}

Una censura positiva indica che le metriche di viaggio sono calcolate su un
sottoinsieme. Non equivale automaticamente a un errore, ma rende il confronto
più delicato e deve essere riportata. La campagna Manhattan usa
\texttt{max\_steps=0}, cioè nessun cutoff imposto dalla pipeline, e termina con
$C=0$ in tutti i 135 casi.

\section{Criteri di confronto}

Il criterio primario è la minore attesa media, a condizione che la censura sia
comparabile e preferibilmente nulla. Si verificano poi tempo di viaggio, time
loss, percentile 95 e velocità. Una variante è considerata sostanzialmente
migliore del base se il vantaggio è coerente sui seed e non deriva da un
numero inferiore di viaggi completati.

Le attivazioni interne vengono usate come evidenza causale debole ma utile. Ad
esempio, risultati e contatori identici tra spillback e base indicano che la
soglia non è stata raggiunta. Non si conclude che lo spillback sia inutile in
generale; si conclude che la campagna non ha testato il regime per cui è stato
progettato.

\section{Ambiente hardware e software}

La campagna finale Manhattan è stata eseguita sul server \texttt{pccabri2} del
dominio \texttt{mat.unimo.it}. L'ambiente comprendeva Ubuntu 22.04.5 LTS,
kernel Linux 5.15, processore AMD Athlon II X4 620 con quattro core, circa
11 GiB di RAM, Python 3.10.12 e binario SUMO 1.27.0. L'ambiente virtuale
conteneva TraCI e sumolib 1.20.0, PyYAML 6.0.1 e NumPy 1.26.4.

Sono stati usati tre worker paralleli. I 135 casi hanno richiesto circa 7 ore e
20 minuti di tempo reale. La GPU non è coinvolta: SUMO e il controller sono
principalmente CPU-bound, e il parallelismo utile deriva dall'esecuzione di più
processi indipendenti.

\section{Minacce alla validità}

\textbf{Validità interna.} Gli scenari condividono popolazioni e codice di
supporto, ma parametri differenti possono modificare più di un meccanismo. Il
caso Nmin, che riduce anche il minimo verde, viene perciò interpretato con i
contatori. Errori del simulatore o della mappa sono mitigati da preflight, log,
test automatici e controllo degli errori SUMO.

\textbf{Validità di costrutto.} Attesa, viaggio e time loss rappresentano
aspetti diversi dell'efficienza, ma non descrivono sicurezza, priorità del
trasporto pubblico o esperienza pedonale. Le conclusioni sono limitate alla
mobilità veicolare. Indicatori ambientali eventualmente presenti negli output
generici non appartengono al modello di valutazione.

\textbf{Validità esterna.} Manhattan è regolare, completamente semaforizzata e
usa domanda balanced. I risultati non si trasferiscono automaticamente a una
rete reale irregolare o a traffico fortemente concentrato. La seconda parte
della tesi affronta questa limitazione con Bologna e con i set skewed e peak.

\textbf{Validità statistica.} Cinque seed permettono di osservare stabilità e
costruire confronti appaiati, ma non sono sufficienti per caratterizzare ogni
coda della distribuzione. Differenze molto piccole tra varianti devono quindi
essere considerate indicative, non definitive. Al contrario, riduzioni di
decine di secondi rispetto al fixed risultano molto maggiori della variabilità
tra seed.

\chapter{Valutazione su Manhattan 8x8}
\label{chap:manhattan}

La rete Manhattan costituisce il primo banco di prova completo perché riduce
le peculiarità topologiche e consente di osservare il comportamento dei
controller in una griglia ripetibile. Il capitolo presenta configurazione,
risultati e contatori della run finale. L'obiettivo non è scegliere un unico
parametro universale, ma capire quali meccanismi producano un beneficio prima
di passare alla rete reale.

\section{Motivazione dello scenario sintetico}

Una rete reale combina contemporaneamente segmenti corti, svolte, incroci con
fasi diverse e domanda spazialmente non uniforme. Se una variante peggiora,
può essere difficile stabilire quale caratteristica abbia causato il risultato.
Manhattan $8\times8$ offre invece incroci quasi omogenei, elevata
semaforizzazione e molte route. È quindi adatta a una ablation controllata.

La regolarità non rende il test banale. Sessanta semafori prendono decisioni
locali e le code si propagano attraverso la griglia. Il piano fisso, inoltre,
non usa offset coordinati: tutti i programmi partono con offset zero. La
campagna misura dunque quanto l'adattamento locale riesca a sfruttare una
domanda diffusa rispetto alla ripetizione dello stesso ciclo.

\section{Topologia della rete}

La mappa \texttt{manhattan8x8\_100pc} contiene 64 giunzioni non interne, delle
quali 60 semaforizzate e quattro angoli regolati a priorità. I 224 archi non
interni hanno complessivamente 448 corsie, due per arco. Le lunghezze delle
corsie sono comprese tra 139,2 e 143,2 metri, con media circa 139,5 metri. Il
file delle route contiene 800 percorsi, da 1 a 15 archi e con lunghezza media
pari a 7,276 archi.

La \cref{fig:manhattan-grid} rappresenta schematicamente la griglia. I nodi
pieni indicano intersezioni semaforizzate e i quattro quadrati bianchi gli
angoli a priorità.

\begin{figure}[htbp]
    \centering
    \begin{tikzpicture}[x=0.72cm,y=0.72cm]
        \foreach \x in {0,...,7}{
            \draw[black!45] (\x,0) -- (\x,7);
        }
        \foreach \y in {0,...,7}{
            \draw[black!45] (0,\y) -- (7,\y);
        }
        \foreach \x in {0,...,7}{
            \foreach \y in {0,...,7}{
                \fill[black!75] (\x,\y) circle (1.25pt);
            }
        }
        \draw[black,fill=white] (-0.10,-0.10) rectangle (0.10,0.10);
        \draw[black,fill=white] (-0.10,6.90) rectangle (0.10,7.10);
        \draw[black,fill=white] (6.90,-0.10) rectangle (7.10,0.10);
        \draw[black,fill=white] (6.90,6.90) rectangle (7.10,7.10);
        \draw[<->] (0,-0.55) -- node[fill=white,inner sep=1pt]{7 isolati regolari} (7,-0.55);
    \end{tikzpicture}
    \caption{Schema della rete Manhattan $8\times8$.}
    \label{fig:manhattan-grid}
\end{figure}

Ogni semaforo possiede due logiche. Il program0, usato dal fixed, ha due fasi
verdi da 42 secondi. Ciascun passaggio comprende 3 secondi di giallo e 3
secondi di all-red, per un ciclo complessivo di 96 secondi. La logica dinamica
mantiene le stesse transizioni ma assegna alle fasi principali una durata
sufficientemente lunga da lasciare al controller esterno la decisione del
cambio. In entrambi i casi le transizioni di sicurezza sono definite nella
mappa e preservate.

\section{Configurazione sperimentale}

La run finale è identificata da
\texttt{run\_0023\_20260612\_131724}. Comprende una mappa, una distribuzione
balanced, tre domande, cinque seed e nove scenari:

\begin{equation}
    1\times1\times3\times5\times9=135\text{ simulazioni}.
\end{equation}

Le partenze sono distribuite uniformemente nei primi 3600 secondi. Low contiene
2400 veicoli, medium 4000 e high 5600. Non è imposto un numero massimo di passi:
ciascun caso prosegue fino a esaurimento. Tutti i casi sono terminati con
successo; in ciascuno il numero di viaggi completati coincide con quello
pianificato e la censura è zero. Non sono stati rilevati nei log avvisi di
collisione, teletrasporto o frenata d'emergenza.

La \cref{tab:manhattan-scenarios} riassume gli scenari. I valori non sono tutte
le combinazioni provate durante lo sviluppo, ma una configurazione definitiva
per ogni famiglia.

\begin{table}[htbp]
    \centering
    \caption{Scenari della campagna Manhattan.}
    \label{tab:manhattan-scenarios}
    \begin{tabularx}{\textwidth}{lX}
        \toprule
        Scenario & Configurazione caratterizzante \\
        \midrule
        \texttt{fixed\_program0} & Programma statico 0 della mappa \\
        \texttt{mp\_base} & Minimo 10 s, massimo 120 s, margine nullo \\
        \texttt{mp\_lta\_g060\_sf050} & $\gamma=0{,}60$, $s=0{,}50$ veicoli/s \\
        \texttt{mp\_nmin\_a080\_f2} & $\alpha=0{,}80$, floor 2, minimo 4 s, $g=0{,}25$ \\
        \texttt{mp\_downstream\_b04} & $\beta=0{,}4$, EMA $\eta=0{,}30$ \\
        \texttt{mp\_spillback\_on90\_off80} & Soglie 0,90/0,80, almeno 3 veicoli fermi \\
        \texttt{mp\_fair\_mu5\_w20} & $\mu=5$, $w_{1/2}=20$ s \\
        \texttt{mp\_platoon\_x4} & Headway e gap 2,1 s, estensione massima 4 s \\
        \texttt{mp\_switch\_rel005} & Margine relativo 5\% \\
        \bottomrule
    \end{tabularx}
\end{table}

\section{Confronto tra Max-Pressure e controllo fisso}

La \cref{tab:wait-all-manhattan} riporta il tempo medio di attesa e la
deviazione standard tra seed. Tutte le varianti Max-Pressure riducono
nettamente l'attesa rispetto al fixed in ciascun livello. Anche la variante
meno favorevole rimane molto distante dalla baseline.

\begin{table}[htbp]
    \centering
    \caption{Tempo medio di attesa su Manhattan, media $\pm$ deviazione standard
    sui cinque seed, in secondi.}
    \label{tab:wait-all-manhattan}
    \resizebox{\textwidth}{!}{%
    \begin{tabular}{lccc}
        \toprule
        Scenario & Low & Medium & High \\
        \midrule
        \texttt{fixed\_program0}          & $65{,}85\pm1{,}11$ & $66{,}89\pm0{,}24$ & $67{,}75\pm0{,}46$ \\
        \texttt{mp\_base}                & $26{,}69\pm0{,}27$ & $29{,}48\pm0{,}35$ & $32{,}07\pm0{,}33$ \\
        \texttt{mp\_lta\_g060\_sf050}   & $30{,}56\pm0{,}46$ & $31{,}45\pm0{,}57$ & $32{,}77\pm0{,}45$ \\
        \texttt{mp\_nmin\_a080\_f2}     & $\mathbf{23{,}54\pm0{,}41}$ & $\mathbf{26{,}70\pm0{,}15}$ & $32{,}03\pm0{,}76$ \\
        \texttt{mp\_downstream\_b04}    & $31{,}76\pm0{,}41$ & $33{,}05\pm0{,}66$ & $34{,}71\pm0{,}61$ \\
        \texttt{mp\_spillback\_on90\_off80} & $26{,}69\pm0{,}27$ & $29{,}48\pm0{,}35$ & $32{,}07\pm0{,}33$ \\
        \texttt{mp\_fair\_mu5\_w20}    & $26{,}60\pm0{,}31$ & $29{,}63\pm0{,}35$ & $\mathbf{31{,}81\pm0{,}10}$ \\
        \texttt{mp\_platoon\_x4}        & $26{,}92\pm0{,}30$ & $29{,}43\pm0{,}49$ & $31{,}94\pm0{,}33$ \\
        \texttt{mp\_switch\_rel005}     & $26{,}69\pm0{,}27$ & $29{,}48\pm0{,}35$ & $32{,}07\pm0{,}33$ \\
        \bottomrule
    \end{tabular}}
\end{table}

Il vincitore low e medium è Nmin; in high il valore minimo appartiene alla
fairness. Rispetto al fixed, Nmin riduce l'attesa del 64,26\% in low e del
60,09\% in medium. In high la fairness ottiene una riduzione del 53,05\%. La
\cref{fig:wait-manhattan} mostra che il divario rispetto al fixed rimane ampio
al crescere della domanda, mentre le migliori varianti MP convergono verso
valori simili.

\begin{figure}[htbp]
    \centering
    \begin{tikzpicture}
    \begin{axis}[
        width=0.90\textwidth,
        height=7.3cm,
        xlabel={Livello di domanda},
        ylabel={Attesa media [s]},
        symbolic x coords={Low,Medium,High},
        xtick=data,
        ymin=20,ymax=72,
        grid=major,
        legend style={at={(0.5,-0.22)},anchor=north,legend columns=2},
        line width=1pt,
        mark size=2.4pt
    ]
        \addplot+[mark=square*] coordinates {(Low,65.849) (Medium,66.895) (High,67.753)};
        \addlegendentry{Fixed program0}
        \addplot+[mark=triangle*] coordinates {(Low,26.695) (Medium,29.482) (High,32.075)};
        \addlegendentry{MP base}
        \addplot+[mark=*] coordinates {(Low,23.536) (Medium,26.699) (High,32.033)};
        \addlegendentry{MP Nmin}
        \addplot+[mark=diamond*] coordinates {(Low,26.599) (Medium,29.629) (High,31.813)};
        \addlegendentry{MP fairness}
    \end{axis}
    \end{tikzpicture}
    \caption{Andamento del tempo medio di attesa per controller selezionati.}
    \label{fig:wait-manhattan}
\end{figure}

Il vantaggio non riguarda soltanto la media. In low il 95-esimo percentile
dell'attesa passa da 137,22 secondi del fixed a 48,00 di Nmin; in medium da
138,80 a 54,20. In high la fairness ottiene 66,00 secondi contro 141,00 del
fixed. La riduzione interessa quindi anche i viaggi maggiormente penalizzati.

\begin{table}[htbp]
    \centering
    \caption{Confronto tra vincitore e fixed per livello di domanda.}
    \label{tab:winners-manhattan}
    \resizebox{\textwidth}{!}{%
    \begin{tabular}{llrrrrrr}
        \toprule
        Domanda & Vincitore & Attesa & $\Delta$ attesa & Viaggio & $\Delta$ viaggio & Time loss & Velocità \\
        & & [s] & [\%] & [s] & [\%] & [s] & [m/s] \\
        \midrule
        Low    & Nmin     & 23,54 & $-64{,}26$ & 137,47 & $-22{,}04$ & 54,58 & 8,33 \\
        Medium & Nmin     & 26,70 & $-60{,}09$ & 144,78 & $-19{,}81$ & 61,57 & 7,94 \\
        High   & Fairness & 31,81 & $-53{,}05$ & 152,23 & $-17{,}44$ & 68,91 & 7,58 \\
        \midrule
        Low fixed    & -- & 65,85 & -- & 176,33 & -- & 93,42  & 6,74 \\
        Medium fixed & -- & 66,89 & -- & 180,56 & -- & 97,32  & 6,58 \\
        High fixed   & -- & 67,75 & -- & 184,39 & -- & 101,06 & 6,43 \\
        \bottomrule
    \end{tabular}}
\end{table}

Le differenze appaiate confermano che il confronto con il fixed è robusto. Il
risparmio medio di attesa è 42,31 secondi in low per Nmin, con intervallo al
95\% $[41{,}16;43{,}47]$; 40,20 secondi in medium, intervallo
$[39{,}92;40{,}47]$; 35,94 secondi in high per fairness, intervallo
$[35{,}30;36{,}58]$. La variabilità tra seed è dunque piccola rispetto al
guadagno.

\section{Valutazione delle singole estensioni}

Il confronto con il fixed dimostra il valore del controllo adattivo nella
griglia, ma non implica che ogni estensione migliori \texttt{mp\_base}. Le
sezioni seguenti separano i due livelli di confronto.

\subsection{Lost-Time-Aware}

LTA riduce il numero di switch decisi dal punteggio: da 6071 a 5273 in low, da
8591 a 7566 in medium e da 10184 a 9106 in high, come media sui seed. Tuttavia
l'attesa aumenta rispettivamente di 3,86, 1,97 e 0,70 secondi rispetto al base.
Gli interventi del massimo verde crescono fortemente in low, da 91 a circa 260.

Il risultato indica che il margine compensa troppo aggressivamente il costo di
transizione nella domanda leggera e ritarda cambi utili. Con traffico elevato
la penalizzazione si riduce, perché le pressioni più grandi superano più
facilmente il margine. LTA non è quindi errata in principio, ma la sua
configurazione non è vantaggiosa nella griglia balanced studiata.

\subsection{Nmin dinamico}

Nmin produce il miglior risultato low e medium. Rispetto a \texttt{mp\_base}
risparmia 3,16 secondi in low, intervallo appaiato al 95\%
$[2{,}55;3{,}77]$, e 2,78 secondi in medium, intervallo
$[2{,}50;3{,}07]$. In high il vantaggio è soltanto 0,04 secondi, con intervallo
$[-0{,}68;0{,}76]$, quindi non distinguibile dalla variabilità osservata.

I passi di hold Nmin sono in media 0 in low, 2,4 in medium e 23,2 in high. Il
forte vantaggio a basso traffico non è quindi causato dalla permanenza dinamica,
che non si attiva, ma soprattutto dal minimo verde ridotto a 4 secondi. Questa
configurazione consente di correggere più rapidamente una fase che ha già
esaurito la domanda. In high l'hold entra in funzione, ma il tempo di viaggio
sale a 155,28 secondi contro 152,59 del base: proteggere il servizio della fase
corrente può ridurre l'attesa media senza migliorare uniformemente l'intero
viaggio.

\subsection{Fairness}

La fairness è quasi neutra in low e medium e nominalmente migliore in high. I
bonus positivi sono applicati in media 5584, 7951 e 9538 volte nei tre livelli;
il meccanismo è quindi attivo. All'aumentare del carico, più fasi conservano
domanda mentre non sono servite e il bonus può correggere squilibri persistenti.

In high il vantaggio sul base è 0,26 secondi. L'intervallo appaiato al 95\% del
risparmio è $[-0{,}08;0{,}60]$ e include zero. Inoltre, considerando il vincitore
di ogni singolo seed tra le varianti di testa, fairness vince due seed, Nmin
due e platoon uno. È corretto definirla vincitrice nominale aggregata, non una
superiorità statistica definitiva.

\subsection{Penalità downstream}

La penalità downstream è la variante peggiore tra gli MP in tutti i livelli.
Rispetto al base aumenta l'attesa di 5,06 secondi in low, 3,57 in medium e 2,64
in high. Il numero di decisioni di switch cresce, mentre gli interventi del
massimo verde quasi scompaiono.

La rete non raggiunge un regime in cui le uscite siano prossime al blocco. La
penalità continua finisce quindi per reagire a normali livelli di occupazione,
modificando frequentemente la classifica delle fasi senza prevenire uno
spillback reale. Il costo relativo diminuisce in high, ma rimane osservabile
anche su viaggio, time loss e velocità.

\subsection{Anti-spillback}

Lo scenario anti-spillback produce valori identici al base fino alla precisione
dei report. Tutti i contatori di blocco, rilascio e passi bloccati sono zero per
ogni domanda. La spiegazione non è che il codice non venga eseguito: nessuna
corsia raggiunge contemporaneamente il punteggio 0,90 e almeno tre veicoli
fermi richiesti per l'attivazione.

Il risultato delimita la campagna. Anche 5600 veicoli balanced costituiscono un
carico elevato per il confronto, ma non una condizione di spillback. Per
valutare la protezione sono necessarie route più concentrate, picchi temporali,
corsie riceventi corte o una domanda maggiore. La successiva introduzione dei
set skewed e peak nasce precisamente da questa osservazione.

\subsection{Platoon extension}

I prolungamenti per plotone aumentano da 113,6 passi medi in low a 227 in
medium e 288,8 in high. Il meccanismo rileva quindi arrivi compatti. Le metriche
restano però molto vicine al base: leggermente peggiori in low, migliori di
0,06 secondi in medium e di 0,13 in high.

L'estensione massima di 4 secondi è prudente e il traffico balanced non produce
corridoi dominanti. Inoltre, nello scenario standalone non è attivo un guard
downstream effettivo. Il test mostra che il riconoscimento dei plotoni funziona,
ma non dimostra un vantaggio sostanziale nella griglia uniforme.

\subsection{Margini di switch}

Il margine relativo del 5\% produce risultati e contatori esattamente uguali al
base. In tutti i casi in cui una fase alternativa supera quella corrente, il
divario è già maggiore del margine, oppure il cambio non sarebbe avvenuto
nemmeno con margine nullo. La banda morta non modifica quindi alcuna decisione
nella campagna.

Questo esito è informativo: la soglia non va considerata stabilizzante soltanto
perché è stata abilitata. Per avere effetto dovrebbe essere aumentata o testata
in uno scenario con pressioni più vicine e oscillanti; un valore più alto,
tuttavia, rischierebbe di riprodurre il problema osservato con LTA.

\section{Analisi dei contatori di attivazione}

La \cref{tab:activation-manhattan} sintetizza il contatore specifico di ciascuna
variante. Le unità sono passi o eventi medi per simulazione; per il base viene
riportata la coppia switch da margine/interventi da massimo verde.

\begin{table}[htbp]
    \centering
    \caption{Contatori medi di attivazione per run.}
    \label{tab:activation-manhattan}
    \begin{tabular}{lrrr}
        \toprule
        Meccanismo & Low & Medium & High \\
        \midrule
        Base: switch / max-green & 6071 / 91 & 8591 / 40 & 10184 / 21 \\
        LTA: switch / max-green & 5273 / 260 & 7566 / 158 & 9106 / 90 \\
        Nmin: passi di hold & 0,0 & 2,4 & 23,2 \\
        Fairness: bonus positivi & 5584 & 7951 & 9538 \\
        Spillback: passi bloccati & 0,0 & 0,0 & 0,0 \\
        Platoon: passi estesi & 113,6 & 227,0 & 288,8 \\
        \bottomrule
    \end{tabular}
\end{table}

Due conclusioni emergono. Primo, ``attivo'' non significa ``utile'': fairness e
platoon intervengono spesso, ma il vantaggio sul base è piccolo; downstream
modifica molte decisioni e peggiora. Secondo, valori finali uguali possono
essere spiegati: spillback e margine relativo non alterano il percorso
decisionale, mentre la coincidenza non è dovuta a un problema di aggregazione.

\section{Confronto complessivo}

La campagna sostiene con chiarezza l'ipotesi principale: su Manhattan
$8\times8$ balanced, il controllo Max-Pressure riduce in modo ampio attesa,
viaggio e time loss rispetto al program0 nativo, mantenendo censura nulla. Il
Max-Pressure base da solo riduce l'attesa del 59,46\%, 55,93\% e 52,66\% nei
tre livelli. Non è quindi necessario un modulo specialistico per ottenere il
beneficio fondamentale.

Tra le estensioni, Nmin è la scelta più efficace a domanda bassa e media,
principalmente grazie al minimo verde più corto. A domanda alta fairness,
platoon, Nmin e base sono raccolti entro 0,26 secondi di attesa. La classifica
nominale non deve oscurare tale prossimità. LTA e downstream mostrano invece
che aggiungere prudenza o informazione non produce automaticamente una politica
migliore: una regola utile in congestione può costare capacità quando il
fenomeno bersaglio è assente.

Il piano fixed rimane quasi costante tra 65,85 e 67,75 secondi di attesa. Questo
non significa che il traffico non aumenti: viaggio e time loss peggiorano e la
velocità scende. Significa che il ciclo statico impone una componente dominante
di ritardo anche in low, mentre il Max-Pressure evita verdi assegnati a
movimenti con poca domanda. La rete regolare e l'assenza di offset coordinati
favoriscono questo vantaggio locale. Non si può estendere la conclusione a un
piano fisso ottimizzato per ogni profilo di domanda.

\section{Indicazioni ricavate per lo scenario reale}

La valutazione Manhattan produce cinque indicazioni per Bologna:

\begin{enumerate}
    \item mantenere \texttt{mp\_base} come riferimento, perché è già robusto e
    ogni specialista deve giustificare il proprio costo;
    \item ridurre il minimo verde può essere molto efficace a basso carico, ma
    va separato concettualmente dall'hold dinamico Nmin;
    \item usare i contatori per verificare l'attivazione: senza congestione
    downstream, lo spillback non può essere giudicato;
    \item introdurre domanda spazialmente e temporalmente non uniforme per
    creare condizioni più realistiche di corridoio e picco;
    \item evitare l'idea di una variante universalmente migliore: la quasi
    parità in high suggerisce di selezionare la modalità in funzione del regime
    locale invece di mantenere sempre la stessa estensione.
\end{enumerate}

Queste indicazioni motivano il passaggio dalla griglia regolare alla rete di
Bologna. In quella mappa il programma statico è più competitivo, la topologia è
eterogenea e il costo degli switch assume un peso maggiore. Il problema non è
più soltanto dimostrare che il Max-Pressure possa funzionare, ma stabilire
quando convenga comportarsi in modo adattivo, quando conservare il piano nativo
e quale specialista attivare.

\ifcapitolifuturi

\chapter{Adattamento alla rete di Bologna}
\label{chap:bologna-progetto}

\section{Caratteristiche della rete}
\tesischeletro

\section{Problemi riscontrati}
\tesischeletro

\subsection{Complessita' della topologia}
\tesischeletro

\subsection{Stalli e congestione}
\tesischeletro

\subsection{Durata delle simulazioni e censura}
\tesischeletro

\section{Preparazione e correzione della mappa}
\tesischeletro

\section{Analisi del programma semaforico program0}
\tesischeletro

\section{Limiti del Max-Pressure puro su Bologna}
\tesischeletro

\subsection{Costo degli switch}
\tesischeletro

\subsection{Prestazioni nel traffico leggero}
\tesischeletro

\subsection{Assenza di coordinamento globale}
\tesischeletro

\section{Calibrazione della domanda}
\tesischeletro

\section{Introduzione delle distribuzioni skewed e peak}
\tesischeletro

\section{Controller ibrido mp\_program0}
\tesischeletro

\subsection{Misura del carico locale}
\tesischeletro

\subsection{Passaggio tra program0 e Max-Pressure}
\tesischeletro

\subsection{Isteresi e stabilita' della selezione}
\tesischeletro

\section{Meta-controller mp\_super}
\tesischeletro

\subsection{Motivazione}
\tesischeletro

\subsection{Load}
\tesischeletro

\subsection{Imbalance}
\tesischeletro

\subsection{Wait imbalance}
\tesischeletro

\subsection{Burstiness}
\tesischeletro

\subsection{Platoon score}
\tesischeletro

\subsection{Saturazione downstream}
\tesischeletro

\subsection{Regole di selezione delle modalita'}
\tesischeletro

\subsection{Isteresi temporale}
\tesischeletro

\section{Evoluzione delle soglie di mp\_super}
\tesischeletro


\chapter{Valutazione finale su Bologna}
\label{chap:bologna-risultati}

\section{Protocollo sperimentale}
\tesischeletro

\section{Risultati con domanda balanced}
\tesischeletro

\subsection{Traffico low}
\tesischeletro

\subsection{Traffico medium}
\tesischeletro

\subsection{Traffico high}
\tesischeletro

\section{Risultati con domanda skewed}
\tesischeletro

\subsection{Traffico low}
\tesischeletro

\subsection{Traffico medium}
\tesischeletro

\subsection{Traffico high}
\tesischeletro

\section{Risultati con domanda peak}
\tesischeletro

\subsection{Traffico low}
\tesischeletro

\subsection{Traffico medium}
\tesischeletro

\subsection{Traffico high}
\tesischeletro

\section{Confronto con fixed program0}
\tesischeletro

\section{Confronto con gli specialisti Max-Pressure}
\tesischeletro

\section{Utilizzo delle modalita' di mp\_super}
\tesischeletro

\section{Analisi dei casi ad alto traffico}
\tesischeletro

\subsection{Relazione tra attesa, throughput e censura}
\tesischeletro

\subsection{Caso balanced high}
\tesischeletro

\subsection{Caso skewed high}
\tesischeletro

\subsection{Caso peak high}
\tesischeletro

\section{Discussione complessiva}
\tesischeletro

\section{Limiti della valutazione}
\tesischeletro


\chapter{Conclusioni}
\label{chap:conclusioni}

\section{Sintesi del lavoro}
\tesischeletro

\section{Risultati principali}
\tesischeletro

\section{Risposta agli obiettivi della tesi}
\tesischeletro

\section{Limiti}
\tesischeletro

\section{Sviluppi futuri}
\tesischeletro


\appendix

\chapter{Configurazioni sperimentali}
\label{app:configurazioni}

\section{Scenari e parametri}
\tesischeletro

\section{Comandi di esecuzione}
\tesischeletro


\chapter{Risultati completi}
\label{app:risultati}

\section{Tabelle complete di Manhattan}
\tesischeletro

\section{Tabelle complete di Bologna}
\tesischeletro


\chapter{Struttura del software}
\label{app:software}

\section{Organizzazione della repository}
\tesischeletro

\section{File principali}
\tesischeletro


\fi

\printbibliography[heading=bibintoc,title={Riferimenti bibliografici}]

\end{document}
